% Introduction - v0 draft.
Agent systems built on large language models (LLMs) are increasingly deployed as personal assistants over a fixed catalogue of skills (calendar, email, messaging, file management) where the same user intents recur many times per day. Recent position work argues that small language models (SLMs) in the 1 to 7 billion parameter range are not merely cheaper but \emph{architecturally better suited} to such narrow, repetitive invocations~\cite{belcak2025slm}, and edge-deployed SLMs have begun to match strong cloud models on function-calling benchmarks~\cite{erdogan2024tinyagent}. Two practical problems remain. First, off-the-shelf SLMs in the 3B class still fail at structured tool use when given full SKILL.md context. Second, full-context ReAct~\cite{yao2022react} agents impose latency and cost penalties even on intents the system has handled successfully many times before.

We present \textbf{A.L.F.R.E.D.} (\textit{Adaptive Local-First Routing and Execution Distillation}), a middleware layer that compiles successful agent trajectories into deterministic \emph{patterns} (typed JSON records carrying a trigger regex, slot schema, command template, and validator) and routes user intents across two tiers: (i) a small-LLM \textit{Reflexer} given only the matched pattern's compressed slice, and (ii) a full-LLM \textit{Thinker} fallback for unmatched intents. The router is feature-engineered (regex, Laplace-smoothed success rates, and a first-window position bonus), so the routing decision itself costs no LLM call.

Our contributions are: (1) a deployable two-tier execution architecture for personal-assistant agents; (2) an empirical evaluation on a 53-task benchmark with Ministral-3-3B-Instruct, reporting 96.2\% router accuracy and 91.4\% atomic gold-command match while preserving roughly 95\% of the larger model's success rate; (3) three named deployment findings, namely success-rate stats poisoning, sibling-pattern hijacking, and a 9.4 percentage-point gain from position-weighted trigger scoring, each grounded in an A/B experiment; and (4) an open-source reference implementation including the benchmark harness.
